\Needspace{6\baselineskip}
\Article{326}
Das vorangehende Verfahren ist schon für sich so bequem, dass es kaum etwas zu wünschen übrig lässt. Dennoch kann es durch verschiedene Hilfsmittel noch erheblich verkürzt werden, von denen wir hier nur wenige berühren können. Wir beschränken die Untersuchung daher auf den Fall, dass die Ausschlusszahl eine ungerade Primzahl ist, die \(A\) nicht teilt, oder eine Potenz einer solchen Primzahl, zumal sich die übrigen Fälle entweder darauf zurückführen oder durch ein ähnliches Verfahren behandeln lassen. Nehmen wir \textit{erstens} an, dass die Ausschlusszahl \(E=p\) eine Primzahl ist, die keine der Zahlen \(m,\allowbreak n\) teilt, und die Werte der Ausdrücke \(\frac Am,\allowbreak -\frac{na}m,\allowbreak -\frac{nb}m\), \(-\frac{nc}m\) usw. \((\mathrm{mod.}\ p)\) der Reihe nach \(k,\allowbreak \mathfrak A,\allowbreak \mathfrak B,\allowbreak \mathfrak C\) usw. sind. Dann findet man die Zahlen \(\alpha,\allowbreak \beta,\allowbreak \gamma\) usw. durch die Kongruenzen \(\alpha\equiv k+\mathfrak A,\allowbreak \ \beta\equiv k+\mathfrak B,\allowbreak \ \gamma\equiv k+\mathfrak C\) usw. \((\mathrm{mod.}\ p)\). Die Zahlen \(\mathfrak A,\allowbreak \mathfrak B,\allowbreak \mathfrak C\) usw. lassen sich jedoch durch ein Hilfsmittel, das dem in Artikel \(322\) verwendeten völlig entspricht, ohne Berechnung der Kongruenzen ermitteln. Sie stimmen entweder mit sämtlichen Nichtresten oder mit sämtlichen Resten von \(p\), außer \(0\), überein, je nachdem, ob der Wert des Ausdrucks \(-\frac mn\ (\mathrm{mod.}\ p)\) oder, was hier auf dasselbe hinausläuft, die Zahl \(-mn\) ein Rest oder ein Nichtrest von \(p\) ist. So hat man im \textit{Beispiel} II des vorigen Artikels für \(E=17\) den Wert \(k=7\); \(-mn=-1365\equiv12\) ist ein Nichtrest von \(17\). Daher sind die Zahlen \(\mathfrak A,\allowbreak \mathfrak B\) usw. gleich \(1,\allowbreak 2,\allowbreak 4,\allowbreak 8,\allowbreak 9,\allowbreak 13,\allowbreak 15,\allowbreak 16\) und die Zahlen \(\alpha,\allowbreak \beta\) usw. gleich \(8,\allowbreak 9,\allowbreak 11,\allowbreak 15,\allowbreak 16,\allowbreak 3,\allowbreak 5,\allowbreak 6\). Davon sind \(8,\allowbreak 9,\allowbreak 15,\allowbreak 16\) Reste; somit ergeben sich für \(\pm h,\allowbreak h'\) usw. die Werte \(\pm5,\allowbreak 3,\allowbreak 7,\allowbreak 4\). — Wer öfter solche Probleme zu lösen hat, wird sich die Arbeit erheblich erleichtern, wenn er für mehrere Primzahlen \(p\) die Werte von \(h,\allowbreak h'\) usw., die den einzelnen Werten von \(k\ (1,\allowbreak 2,\allowbreak 3\ldots p-1)\) entsprechen, unter beiden Voraussetzungen berechnet: dass \(-mn\) Rest beziehungsweise Nichtrest von \(p\) ist. Ferner bemerken wir, dass die Anzahl der Zahlen \(h,\allowbreak -h,\allowbreak h'\) usw. stets \(\frac12(p-1)\) ist, wenn sowohl \(k\) als auch \(-mn\) Reste oder beide Nichtreste von \(p\) sind; \(\frac12(p-3)\), wenn die erste Zahl \(R.\) und die zweite \(NR.\) ist; und \(\frac12(p+1)\), wenn die erste \(NR.\) und die zweite \(R.\) ist. Den Beweis dieses Satzes müssen wir jedoch auslassen, um nicht zu ausführlich zu werden.

\textit{Zweitens} lassen sich die Fälle, in denen \(E\) eine Primzahl ist, die \(n\) teilt, oder eine Potenz einer ungeraden Primzahl, die \(n\) teilt oder nicht teilt, noch bequemer behandeln. Wir fassen sämtliche Fälle zusammen und setzen unter Beibehaltung aller Bezeichnungen des Artikels \(324\) nun \(n=n'p^\nu\), sodass \(n'\) durch \(p\) nicht teilbar ist. Die Zahlen \(a,\allowbreak b,\allowbreak c\) usw. sind die Produkte von \(p^{\mu-1}\) entweder mit sämtlichen Zahlen kleiner als \(p\), außer \(0\), oder mit sämtlichen Nichtresten von \(p\) unterhalb \(p\), je nachdem, ob \(\mu\) gerade oder ungerade ist; sie seien allgemein durch \(up^{\mu-1}\) bezeichnet. Es sei \(k\) der Wert des Ausdrucks \(\frac Am\ (\mathrm{mod.}\ p^{\mu+\nu})\). Dieser ist durch \(p\) nicht teilbar, weil dieselbe Eigenschaft bei \(A\) vorausgesetzt wird;
\SourceLeaf{402}
ferner ist klar, dass sämtliche Zahlen \(\alpha,\allowbreak \beta,\allowbreak \gamma\) usw. der Zahl \(k\) nach dem Modul \(p\) kongruent sind und daher \(p^\mu\) nichts aus \(\Omega\) ausschließt, wenn \(kNp\) gilt. Gilt dagegen \(kRp\) und somit auch \(kRp^{\mu+\nu}\), so sei \(r\) ein Wert des Ausdrucks \(\surd k\ (\mathrm{mod.}\ p^{\mu+\nu})\), der durch \(p\) nicht teilbar ist, und \(e\) der Wert von \(-\frac{n'}{2mr}\ (\mathrm{mod.}\ p)\). Dann gilt \(\alpha\equiv rr+2era p^\nu\ (\mathrm{mod.}\ p^{\mu+\nu})\), woraus leicht folgt, dass \(\alpha\) ein Rest von \(p^{\mu+\nu}\) ist und die Werte des Ausdrucks \(\surd\alpha\ (\mathrm{mod.}\ p^{\mu+\nu})\) gleich \(\pm(r+ea p^\nu)\) sind. Somit sind sämtliche Zahlen \(h,\allowbreak h',\allowbreak h''\) usw. durch \(r+uep^{\mu+\nu-1}\) darstellbar. Schließlich folgt daraus leicht, dass die Zahlen \(h,\allowbreak h',\allowbreak h''\) usw. entstehen, indem man \(r\) zu den Produkten von \(p^{\mu+\nu-1}\) entweder mit sämtlichen Zahlen unter \(p\), außer \(0\), addiert, wenn \(\mu\) gerade ist; oder mit sämtlichen Nichtresten von \(p\) unter dieser Grenze, wenn \(\mu\) ungerade ist und \(eRp\) oder, was hier dasselbe bedeutet, \(-2mrn'Rp\) gilt; oder mit sämtlichen Resten, außer \(0\), wenn \(\mu\) ungerade ist und \(-2mrn'Np\) gilt.

Sobald die Zahlen \(h,\allowbreak h'\) usw. für die einzelnen Ausschlusszahlen, die man verwenden will, ermittelt sind, lässt sich im Übrigen auch der Ausschluss selbst durch mechanische Operationen ausführen, die sich jeder in diesen Dingen Erfahrene leicht selbst ausdenken kann, wenn es der Mühe wert erscheint.

Schließlich müssen wir bemerken, dass jede Gleichung \(axx+2bxy+cyy=M\), in der \(bb-ac\) eine negative Zahl \(=-D\) ist, leicht auf die oben betrachtete Form zurückgeführt werden kann. Bezeichnen wir nämlich den größten gemeinsamen Teiler der Zahlen \(a,\allowbreak b\) mit \(m\) und setzen \[ a=ma'.\quad b=mb',\quad \frac Dm=a'c-mb'b'=n,\quad a'x+b'y=x' \] so ist jene Gleichung offenbar zu \(mx'x'+nyy=a'M\) äquivalent und lässt sich nach den oben angegebenen Regeln lösen. Von deren Lösungen sind aber nur diejenigen beizubehalten, in denen \(x'-b'y\) durch \(a'\) teilbar ist, aus denen sich also ganzzahlige Werte von \(x\) ergeben.

\Needspace{9\baselineskip}\begin{center}\textit{Ein anderes Verfahren zur Lösung der Kongruenz \(xx\equiv A\), wenn \(A\) negativ ist.}\end{center}


\Needspace{6\baselineskip}\Article{327}
Wie die in Abschnitt V angegebene direkte Lösung der Gleichung \(axx+2bxy+cyy=M\) die Werte des Ausdrucks \(\surd{(bb-ac)}\ (\mathrm{mod.}\ M)\) als bekannt voraussetzt, so liefert umgekehrt für den Fall, dass \(bb-ac\) negativ ist, die oben erläuterte indirekte Lösung ein sehr bequemes Verfahren zur Ermittlung jener Werte. Besonders bei einem sehr großen Wert von \(M\) ist es dem Verfahren des Artikels \(322\) und der folgenden Artikel bei weitem vorzuziehen. Wir nehmen jedoch an, dass \(M\) eine Primzahl ist oder wenigstens ihre Faktoren, falls sie zusammengesetzt ist,
\SourceLeaf{403}
noch unbekannt sind. Wüsste man nämlich, dass die Primzahl \(p\) die Zahl \(M\) teilt und \(M=p^\mu M'\) gilt, wobei \(M'\) den Faktor \(p\) nicht mehr enthält, so wäre es viel bequemer, die Werte des Ausdrucks \(\surd{(bb-ac)}\) für die Moduln \(p^\mu\) und \(M'\) einzeln zu ermitteln, die ersteren aus den Werten nach dem Modul \(p\), Artikel \(101\), und die Werte nach dem Modul \(M\) durch deren Verbindung abzuleiten, Artikel \(105\).

Gesucht seien also sämtliche Werte des Ausdrucks \(\surd{-D}\ (\mathrm{mod.}\ M)\), wobei \(D\) und \(M\) als positiv und \(M\) unter einer Form der Teiler von \(xx+D\) enthalten vorausgesetzt werden, Artikel \(147\) und die folgenden. Andernfalls stünde nämlich von vornherein fest, dass keine Zahlen dem gegebenen Ausdruck genügen können. Die gesuchten Werte, von denen stets je zwei entgegengesetzt sind, seien \(\pm r,\allowbreak \pm r',\allowbreak \pm r''\) usw., und es gelte \(D+rr=Mh,\allowbreak \ D+r'r'=Mh',\allowbreak \ D+r''r''=Mh''\) usw. Die Klassen, zu denen die Formen \((M,\allowbreak r,\allowbreak h),\allowbreak \ (M,\allowbreak -r,\allowbreak h),\allowbreak \ (M,\allowbreak r',\allowbreak h')\), \((M,\allowbreak -r',\allowbreak h'),\allowbreak \ (M,\allowbreak r'',\allowbreak h'').\ (M,\allowbreak -r'',\allowbreak h'')\) usw. gehören, seien der Reihe nach mit \(\mathfrak C,\allowbreak -\mathfrak C,\allowbreak \mathfrak C',\allowbreak -\mathfrak C'\), \(\mathfrak C'',\allowbreak -\mathfrak C''\) usw. und ihre Gesamtheit mit \(\mathfrak G\) bezeichnet. Diese Klassen sind im Allgemeinen zwar als unbekannt anzusehen. Dennoch ist klar, \textit{erstens}, dass sie sämtlich positiv und eigentlich primitiv sind, \textit{zweitens}, dass sie alle zu demselben Geschlecht gehören, dessen \textit{Charakter} sich leicht aus der Beschaffenheit der Zahl \(M\), \textit{d. h.} aus ihren Beziehungen zu den einzelnen Primteilern von \(D\) und außerdem zu \(4\) oder \(8\), wenn diese nötig sind, erkennen lässt, Artikel \(230\). Da vorausgesetzt ist, dass \(M\) unter einer Form der Teiler von \(xx+D\) enthalten ist, können wir von vornherein sicher sein, dass diesem Charakter notwendig ein Geschlecht positiver, eigentlich primitiver Formen der Determinante \(-D\) entspricht, selbst wenn der Ausdruck \(\surd{-D}\ (\mathrm{mod.}\ M)\) möglicherweise keine Lösungen besitzt. Da dieses Geschlecht somit bekannt ist, lassen sich sämtliche darin enthaltenen Klassen ermitteln; sie seien \(C,\allowbreak C',\allowbreak C''\) usw. und ihre Gesamtheit \(G\).

Somit ist klar, dass jede Klasse \(\mathfrak C,\allowbreak -\mathfrak C\) usw. mit einer Klasse in \(G\) identisch sein muss. Es kann auch vorkommen, dass mehrere Klassen in \(\mathfrak G\) untereinander und somit mit derselben Klasse in \(G\) identisch sind. Enthält \(G\) nur eine einzige Klasse, so stimmen gewiss alle Klassen in \(\mathfrak G\) mit dieser überein. Wählt man daher aus den Klassen \(C,\allowbreak C',\allowbreak C''\) usw. die einfachsten Formen \(f,\allowbreak f',\allowbreak f''\) usw., je eine aus jeder, so befindet sich unter diesen eine Form aus jeder Klasse in \(\mathfrak G\). Ist nun \(axx+2bxy+cyy\) eine Form in der Klasse \(\mathfrak C\), so gibt es durch sie zwei Darstellungen der Zahl \(M\), die zum Wert \(r\) gehören; ist die eine \(x=m,\allowbreak y=n\), so ist die andere \(x=-m,\allowbreak y=-n\). Nur der Fall \(D=1\) ist auszunehmen, in dem es vier Darstellungen gibt; siehe Artikel \(180\).

Hieraus folgt, dass, wenn sämtliche Darstellungen der Zahl \(M\) durch die einzelnen Formen \(f,\allowbreak f',\allowbreak f''\) usw. nach dem oben angegebenen indirekten Verfahren ermittelt und
\SourceLeaf{404}
daraus die Werte des Ausdrucks \(\surd{-D}\ (\mathrm{mod.}\ M)\) abgeleitet werden, zu denen die einzelnen Darstellungen gehören, Artikel 154 und die folgenden, man \textit{sämtliche} Werte dieses Ausdrucks erhält, und zwar jeden zweimal oder, wenn \(D=1\), viermal. \textit{Q. E. F.} Finden sich unter \(f,\allowbreak f'\) usw. Formen, durch die \(M\) nicht darstellbar ist, so zeigt dies an, dass sie zu keiner Klasse in \(\mathfrak G\) gehören und deshalb unberücksichtigt bleiben müssen. Lässt sich aber \(M\) durch keine jener Formen darstellen, so muss \(-D\) ein quadratischer Nichtrest von \(M\) sein. --- Bei diesen Operationen sind noch die folgenden Bemerkungen zu beachten.

I. Unter den hier verwendeten Darstellungen der Zahl \(M\) durch die Formen \(f,\allowbreak f'\) usw. sind solche zu verstehen, bei denen die Werte der Unbestimmten teilerfremd sind. Ergeben sich andere, bei denen diese Werte einen gemeinsamen Teiler \(\mu\) haben, was nur vorkommen kann, wenn \(\mu\mu\) die Zahl \(M\) teilt, und gewiss vorkommt, wenn \(-D\,\mathrm R\,\frac{M}{\mu\mu}\) gilt, so müssen sie für den vorliegenden Zweck ganz unberücksichtigt bleiben, obwohl sie in anderer Hinsicht nützlich sein können.

II. Unter sonst gleichen Umständen ist die Arbeit offenbar umso leichter, je kleiner die Anzahl der Klassen \(f,\allowbreak f',\allowbreak f''\) usw. ist, und daher am kürzesten, wenn \(D\) eine der 65 in Artikel 303 angegebenen Zahlen ist, für die jedes Geschlecht nur eine einzige Klasse enthält.

III. Da je zwei solche Darstellungen \(x=m,\allowbreak y=n\) und \(x=-m,\allowbreak y=-n\) stets zu demselben Wert gehören, genügt es offenbar, nur diejenigen Darstellungen zu betrachten, in denen \(y\) positiv ist. Verschiedene Darstellungen \textit{dieser Art} entsprechen daher stets verschiedenen Werten des Ausdrucks \(\surd{-D}\ (\mathrm{mod.}\ M)\), sodass die Anzahl sämtlicher verschiedener Werte der Anzahl aller sich ergebenden Darstellungen dieser Art gleich ist, stets mit Ausnahme des Falles \(D=1\), in dem die erstere Anzahl die Hälfte der letzteren ist.

\begin{samepage}
IV. Da mit dem einen der beiden entgegengesetzten Werte \(+r,\allowbreak -r\) auch der andere unmittelbar bekannt ist, lassen sich die Operationen noch etwas verkürzen. Ergibt sich der Wert \(r\) aus einer Darstellung der Zahl \(M\) durch eine Form in der Klasse \(C\), \textit{d. h.} gilt \(\mathfrak C=C\), so entsteht der entgegengesetzte Wert \(-r\) offenbar aus einer Darstellung durch eine Form in der zu \(C\) entgegengesetzten Klasse. Diese unterscheidet sich von der Klasse \(C\), sofern letztere nicht ambig ist. Sind also nicht sämtliche Klassen in \(G\) ambig, so braucht man von den übrigen nur die Hälfte zu betrachten, nämlich je eine aus jedem entgegengesetzten Paar. Die andere kann unberücksichtigt bleiben, da man schon ohne Rechnung voraussieht, dass sich aus ihr die entgegengesetzten Werte zu denen der ersteren ergeben. Ist dagegen \(C\) ambig, so entstehen daraus die beiden Werte \(r\) und \(-r\) zugleich. Wurde nämlich aus \(C\) die ambige Form
\SourceLeaf{405}
\(axx+2bxy+cyy\) gewählt und ergab sich der Wert \(r\) aus der Darstellung \(x=m,\allowbreak y=n\), so ergibt sich der Wert \(-r\) aus \(x=-m-\frac{2bn}{a},\allowbreak \ y=n\).
\par\end{samepage}

V. Im Fall \(D=1\) gibt es insgesamt nur eine Klasse, aus der, wie wir annehmen dürfen, die Form \(xx+yy\) gewählt worden sei. Ergibt sich der Wert \(r\) aus der Darstellung \(x=m,\allowbreak y=n\), so ergibt sich derselbe aus \(x=-m,\allowbreak y=-n;\allowbreak \ x=n,\allowbreak y=-m;\allowbreak \ x=-n,\allowbreak y=m\) und der entgegengesetzte Wert \(-r\) aus \(x=m,\allowbreak y=-n;\allowbreak x=-m,\allowbreak y=n;\allowbreak \ x=n,\allowbreak y=m;\allowbreak \ x=-n,\allowbreak y=-m\). Daher genügt von diesen acht Darstellungen, die eine einzige Zerlegung bilden, eine, wenn wir nur dem daraus erhaltenen Wert seinen entgegengesetzten zuordnen.

VI. Der Wert des Ausdrucks \(\surd{-D}\ (\mathrm{mod.}\ M)\), zu dem die Darstellung \(M=amm+2bmn+cnn\) gehört, ist nach Artikel 155 gleich \(\mu(mb+nc)-\nu(ma+nb)\) oder einer beliebigen hierzu nach \(M\) kongruenten Zahl, wobei \(\mu,\allowbreak \nu\) so gewählt sind, dass \(\mu m+\nu n=1\) gilt. Bezeichnen wir einen solchen Wert mit \(v\), so gilt \[\begin{aligned}
mv&\equiv\mu m(mb+nc)-\nu(M-mnb-nnc)\\
&\equiv(\mu m+\nu n)(mb+nc)\equiv mb+nc\ (\mathrm{mod.}\ M)
\end{aligned}\] Somit ist \(v\) ein Wert des Ausdrucks \(\frac{mb+nc}{m}\ (\mathrm{mod.}\ M)\); ebenso findet man, dass \(v\) ein Wert des Ausdrucks \(-\frac{ma+nb}{n}\ (\mathrm{mod.}\ M)\) ist. Diese Formeln sind häufig derjenigen vorzuziehen, aus der sie abgeleitet wurden.

\Needspace{6\baselineskip}\Article{328}
{\rightskip=0pt plus 1fil\relax
\textit{Beispiele.} I. Gesucht sind sämtliche Werte des Ausdrucks \(\surd{-1365}\ (\mathrm{mod.}\ 5428681=M)\). Die Zahl \(M\) ist hier \(\equiv1,\allowbreak 1,\allowbreak 1,\allowbreak 6,\allowbreak 11\ (\mathrm{mod.}\ 4,\allowbreak 3,\allowbreak 5,\allowbreak 7,\allowbreak 13)\) und somit unter einer Form der Teiler von \(xx+1,\allowbreak \ xx+3,\allowbreak \ xx-5\) sowie unter einer Form der Nichtteiler von \(xx+7,\allowbreak \ xx-13\), also unter einer Form der Teiler von \(xx+1365\) enthalten. Der Charakter des Geschlechts, in dem sich die Klassen \(\mathfrak G\) befinden, ist daher \(1,\allowbreak 4;\allowbreak \ R3;\allowbreak \ R5;\allowbreak \ N7;\allowbreak \ N13\). Dieses Geschlecht enthält nur eine Klasse, aus der wir die Form \(6xx+6xy+229yy\) wählen. Um sämtliche Darstellungen der Zahl \(M\) durch diese zu finden, setzen wir \(2x+y=x'\), woraus \(3x'x'+455yy=2M\) folgen muss. Diese Gleichung besitzt vier Lösungen, in denen \(y\) positiv ist, nämlich \(y=127,\allowbreak \ x'=\pm1083,\allowbreak \ y=119,\allowbreak \ x'=\pm1213\). Daraus ergeben sich vier Lösungen der Gleichung \(6xx+6xy+229yy=M\), in denen \(y\) positiv ist: \[\renewcommand{\arraystretch}{1.25}\begin{array}{c|r|r|r|r}x&478&-605&547&-666\\y&127&127&119&119\end{array}\] Die erste Lösung liefert für \(v\) den Wert des Ausdrucks \(\frac{30517}{478}\) oder \(-\frac{3249}{127}\ (\mathrm{mod.}\ M)\), woraus man
\SourceLeaf{406}
\(2350978\) erhält; die zweite liefert den entgegengesetzten Wert \(-2350978\), die dritte \(2600262\) und die vierte den entgegengesetzten Wert \(-2600262\).
\par}

II. Sollen die Werte des Ausdrucks \(\surd{-286}\ (\mathrm{mod.}\ 4272943=M)\) gefunden werden, so ergibt sich als Charakter des Geschlechts, in dem die Klassen \(\mathfrak G\) enthalten sind, \(1\text{ und }7,\allowbreak 8;\allowbreak \ R11;\allowbreak \ R13\). Es ist somit das Hauptgeschlecht, das drei Klassen enthält, dargestellt durch die Formen \((1,\allowbreak 0,\allowbreak 286)\), \((14,\allowbreak 6,\allowbreak 23),\allowbreak \ (14,\allowbreak -6,\allowbreak 23)\). Die dritte darf man als der zweiten entgegengesetzt unberücksichtigt lassen. Durch die Form \(xx+286yy\) findet man zwei Darstellungen der Zahl \(M\), in denen \(y\) positiv ist, nämlich \(y=103,\allowbreak \ x=\pm1113\), woraus sich die Werte \(1493445,\allowbreak \ -1493445\) des gegebenen Ausdrucks ergeben. Durch die Form \((14,\allowbreak 6,\allowbreak 23)\) ist aber \(M\) nicht darstellbar; daraus folgt, dass es außer den beiden gefundenen keine weiteren Werte gibt.

III. Ist der Ausdruck \(\surd{-70}\ (\mathrm{mod.}\ 997331)\) gegeben, so müssen die Klassen \(\mathfrak G\) in dem Geschlecht liegen, dessen Charakter \(3\text{ und }5,\allowbreak 8;\allowbreak \ R5;\allowbreak \ N7\) ist. Dieses enthält nur eine Klasse, deren darstellende Form \((5,\allowbreak 0,\allowbreak 14)\) ist. Führt man jedoch die Rechnung aus, so findet man, dass die Zahl \(997331\) durch die Form \((5,\allowbreak 0,\allowbreak 14)\) nicht darstellbar ist; daher muss \(-70\) ein quadratischer Nichtrest jener Zahl sein.
